% Source: https://artofproblemsolving.com/community/c3h1642458_infinite_series
% Stable identity: aops-topic-1642458-q01
\textbf{(a)} All terms are positive and the first bracket alone is the harmonic series without its first term, which diverges: grouping $\dfrac{1}{3}+\dfrac{1}{4}\ge\dfrac{1}{2}$, $\dfrac{1}{5}+\cdots+\dfrac{1}{8}\ge\dfrac{1}{2}$, $\dfrac{1}{9}+\cdots+\dfrac{1}{16}\ge\dfrac{1}{2}$, and so on, the partial sums exceed any bound. Hence the whole series diverges.

\textbf{(b)} All terms are positive, so we may sum in any order. Summing over $n$ first, for fixed $k\ge2$ we get a geometric series:
\[
\sum_{n=2}^{\infty}\dfrac{1}{k^n}=\dfrac{\dfrac{1}{k^2}}{1-\dfrac{1}{k}}=\dfrac{1}{k(k-1)}=\dfrac{1}{k-1}-\dfrac{1}{k} .
\]
Then the sum over $k$ telescopes:
\[
\sum_{k=2}^{\infty}\left(\dfrac{1}{k-1}-\dfrac{1}{k}\right)=\boxed{1}.
\]
